\documentclass[a4paper]{article}

% Paquetes básicos
% Español (México) con XeLaTeX: fontspec para fuentes Unicode, polyglossia
% para la división silábica y los nombres del documento en la variante mexicana.
\usepackage{fontspec}
\usepackage{polyglossia}
\setdefaultlanguage[variant=mexican]{spanish}
% Los nombres chinos van como «pinyin (original)»: xeCJK da a los caracteres
% Han su propia fuente; sin ella XeLaTeX los omite con un aviso.
% FandolSong no cubre algunos caracteres de nombres (U+73FA); WenQuanYi Zen Hei sí.
\usepackage[AutoFallBack=true]{xeCJK}
\setCJKmainfont{FandolSong-Regular.otf}
\setCJKfallbackfamilyfont{\CJKrmdefault}{WenQuanYi Zen Hei}
% polyglossia no traduce los nombres de listados.
\AtBeginDocument{\providecommand{\lstlistingname}{}\renewcommand{\lstlistingname}{Listado}%
  \providecommand{\lstlistlistingname}{}\renewcommand{\lstlistlistingname}{Índice de listados}}
\usepackage{amsmath, amssymb}
\usepackage{graphicx}
\usepackage[margin=2.5cm]{geometry}
\usepackage[most]{tcolorbox}
\usepackage{etoolbox}
\usepackage{listings}
\usepackage{hyperref}
\usepackage{booktabs}
\usepackage{subcaption}
\usepackage{float}

% Cajas de resaltado
\newtcolorbox{knowledgebox}[1]{
    enhanced, colback=blue!5!white, colframe=blue!75!black, colbacktitle=blue!75!black,
    coltitle=white, fonttitle=\bfseries, title={#1}, attach boxed title to top left={yshift=-2mm, xshift=2mm},
    boxrule=1pt, sharp corners
}

\newtcolorbox{importantbox}[1]{
    enhanced, colback=yellow!10!white, colframe=yellow!80!black, colbacktitle=yellow!80!black,
    coltitle=black, fonttitle=\bfseries, title={#1}, sharp corners
}

\newtcolorbox{warningbox}[1]{
    enhanced, colback=red!5!white, colframe=red!75!black, colbacktitle=red!75!black,
    coltitle=white, fonttitle=\bfseries, title={#1}, sharp corners
}

% Listados
\lstset{
    language=Python,
    basicstyle=\ttfamily\small,
    keywordstyle=\color{blue},
    stringstyle=\color{red!60!black},
    commentstyle=\color{green!60!black},
    breaklines=true,
    frame=single,
    numbers=left,
    numberstyle=\tiny\color{gray},
    captionpos=b,
    extendedchars=true
}

% Metadatos de la portada
\newcommand{\notetitle}{KAIST CS492D Lección 3: Modelos Probabilísticos de Difusión con Reducción de Ruido (Parte 2)}
\newcommand{\noteauthors}{Elaboradas a partir de la clase de Minhyuk Sung}
\newcommand{\notedate}{Otoño de 2025}
\newcommand{\videochannel}{Minhyuk Sung (KAIST)}
\newcommand{\videopublishdate}{2025}
\newcommand{\videoduration}{59 minutos}
\newcommand{\videourl}{https://www.youtube.com/watch?v=5_t3uSvdrWg}
\newcommand{\videocoverpath}{cover.jpg}
\newcommand{\repourl}{https://github.com/hqhq1025/ai-course-notes}

% Pie de página con información del repositorio
\usepackage{fancyhdr}
\pagestyle{fancy}
\fancyhf{}
\fancyfoot[C]{\thepage}
\fancyfoot[R]{\footnotesize\color{gray}\href{\repourl}{AI Course Notes}}
\renewcommand{\headrulewidth}{0pt}
\renewcommand{\footrulewidth}{0pt}

\begin{document}
\begin{titlepage}
\centering
{\Large Notas del curso\par}
\vspace{1.2cm}
{\huge\bfseries \notetitle\par}
\vspace{0.8cm}
{\large \noteauthors\par}
\vspace{0.3cm}
{\large \notedate\par}
\vspace{1.2cm}

\ifdefempty{\videocoverpath}{
{\small Al generar las notas, indique la ruta local de la portada del video.\par}
}{
\includegraphics[width=0.82\textwidth,height=0.45\textheight,keepaspectratio]{\videocoverpath}\par
}

\vfill
\begin{tcolorbox}[width=0.9\textwidth, colback=black!2!white, colframe=black!60, sharp corners]
\textbf{Autor o canal del video}: \videochannel\par
\textbf{Fecha de publicación}: \videopublishdate\par
\textbf{Duración del video}: \videoduration\par
\ifdefempty{\videourl}{}{
\textbf{Enlace del video}: \href{\videourl}{\nolinkurl{\videourl}}\par
}\par
\textbf{Página del proyecto}: \href{\repourl}{\nolinkurl{\repourl}}\par
\end{tcolorbox}
\end{titlepage}

\tableofcontents
\newpage

%% ========== Sección 1: Revisión y motivación ==========


\section{Repaso y motivación}

Esta clase es la segunda parte de la serie sobre DDPM (Denoising Diffusion Probabilistic Models). En la clase anterior, conocimos el marco básico de los modelos de difusión: partiendo de las limitaciones del VAE, introdujimos la idea de variables latentes múltiples y definimos así el proceso hacia adelante y el proceso hacia atrás de los modelos de difusión. En esta clase comenzaremos con la deducción del objetivo de entrenamiento, derivaremos paso a paso una función de pérdida de entrenamiento simplificada e introduciremos tres formas de parametrización equivalentes: predicción de la media, predicción de la muestra limpia y predicción del ruido.

\subsection{Motivación desde el VAE hasta los modelos de difusión}

Repasemos el objetivo central de los modelos generativos: modelar la distribución marginal de datos $p(x_0)$. En el marco del VAE, esto se logra mediante la combinación de una distribución previa $p(z)$ (gausiana estándar) y una distribución de verosimilitud $p_\theta(x|z)$ (decodificador):

$$p_\theta(x_0) = \int p(z) \, p_\theta(x_0 | z) \, dz$$

La limitación del VAE radica en que tanto la previa como la verosimilitud se modelan como distribuciones gaussianas; esto limita la capacidad de expresión de la distribución marginal $p_\theta(x_0)$: es esencialmente una mezcla de distribuciones gaussianas y, por tanto, puede no ser capaz de capturar plenamente las características multimodales de distribuciones de datos complejas.

\begin{importantbox}{Idea central de los modelos de difusión}
Los modelos de difusión aumentan la capacidad de expresión de la distribución de verosimilitud introduciendo una secuencia de variables latentes \textbf{múltiples} $x_1, x_2, \dots, x_T$. Aunque la distribución de transición en cada paso sigue siendo gaussiana, mediante la combinación de transformaciones no lineales múltiples, la distribución de verosimilitud global puede expresar distribuciones mucho más complejas que una única gaussiana.
\end{importantbox}

\subsection{Configuración básica de los modelos de difusión}

Los componentes centrales de los modelos de difusión incluyen:

\begin{itemize}
    \item \textbf{Punto de datos} $x_0$: muestra limpia muestreada de la distribución de datos reales $q(x_0)$
    \item \textbf{Variable latente final} $x_T$: muestra de una distribución gaussiana estándar $\mathcal{N}(0, I)$
    \item \textbf{Variables latentes intermedias} $x_1, \dots, x_{T-1}$: estados intermedios entre los datos y el ruido puro
    \item \textbf{Proceso hacia adelante} (Forward Process): de $x_0$ a $x_T$, inyectando ruido gradualmente
    \item \textbf{Proceso hacia atrás} (Reverse Process): de $x_T$ a $x_0$, eliminando el ruido gradualmente
\end{itemize}

\begin{knowledgebox}{Diferencia clave con el VAE}
En el VAE, tanto el codificador (distribución posterior) como el decodificador (distribución de verosimilitud) deben aprenderse. En cambio, en los modelos de difusión, el proceso hacia adelante está \textbf{predefinido} (no requiere aprendizaje); solo debemos aprender el proceso hacia atrás. Además, todas las variables latentes $x_t$ y los datos $x_0$ tienen la \textbf{misma dimensión}, lo cual difiere del caso habitual en el VAE donde las variables latentes suelen tener menor dimensión.
\end{knowledgebox}

Sobre la propiedad de Markov: tanto el proceso hacia adelante como el hacia atrás se definen como procesos de Markov, es decir, $x_t$ depende únicamente de $x_{t-1}$ (hacia adelante) o de $x_{t+1}$ (hacia atrás), sin depender de estados anteriores. Formalmente:

$$q(x_t | x_{t-1}, x_0) = q(x_t | x_{t-1})$$

Esta propiedad simplifica enormemente las deducciones matemáticas posteriores.

\subsection{Resumen de la sección}

Los modelos de difusión logran una gran capacidad de expresión manteniendo la simplicidad del entrenamiento mediante un proceso predefinido de inyección de ruido hacia adelante y un proceso hacia atrás de eliminación de ruido que se aprende. A diferencia del VAE, los modelos de difusión no requieren aprender un codificador y todas las variables latentes comparten la misma dimensión que los datos.

%% ========== Segunda sección: Proceso hacia adelante ==========


\section{Detalle del proceso de difusión hacia adelante}

\subsection{Definición de la distribución de transición hacia adelante}

Cada paso de transición del proceso de difusión se define como una distribución gaussiana de forma específica:

$$q(x_t | x_{t-1}) = \mathcal{N}\big(x_t;\, \sqrt{1 - \beta_t}\, x_{t-1},\; \beta_t I\big)$$

Los significados de los diversos símbolos son los siguientes:
\begin{itemize}
    \item $\beta_t \in (0, 1)$: parámetro del cronograma de ruido (noise schedule) para el paso $t$, que es una función \textbf{no decreciente} con respecto al tiempo $t$.
    \item $\sqrt{1 - \beta_t}\, x_{t-1}$: media — escalar el valor del paso anterior hacia el origen.
    \item $\beta_t I$: varianza — inyectar ruido gaussiano isótropo.
\end{itemize}

\begin{importantbox}{Comprensión intuitiva del proceso de difusión hacia adelante}
Cada transición hacia adelante en el proceso realiza dos cosas: (1) multiplicar el valor actual $x_{t-1}$ por un coeficiente ligeramente menor que 1, $\sqrt{1-\beta_t}$, para contraerlo hacia el origen; y (2) agregar ruido gaussiano con varianza $\beta_t$. Cuando $\beta_t$ se acerca a 0, $x_t \approx x_{t-1}$ (cambia casi nada); cuando $\beta_t$ se acerca a 1, $x_t$ se vuelve casi puro ruido. A medida que $t$ aumenta, $\beta_t$ crece gradualmente y el ruido se acumula progresivamente, hasta que $x_T$ tiende a una distribución gaussiana estándar.
\end{importantbox}

\subsection{Distribución de salto hacia adelante (Forward Jump Distribution)}

Una ventaja clave del proceso de difusión hacia adelante es que no necesitamos ejecutarlo paso a paso; podemos muestrear directamente $x_t$ en cualquier paso temporal $t$ partiendo de $x_0$. Definamos $\alpha_t = 1 - \beta_t$ y $\bar{\alpha}_t = \prod_{s=1}^{t} \alpha_s$, entonces:

$$q(x_t | x_0) = \mathcal{N}\big(x_t;\, \sqrt{\bar{\alpha}_t}\, x_0,\; (1 - \bar{\alpha}_t) I\big)$$

Explicación de los símbolos:
\begin{itemize}
    \item $\bar{\alpha}_t = \prod_{s=1}^{t}(1-\beta_s)$: coeficiente acumulado conservado, que disminuye con $t$.
    \item $\sqrt{\bar{\alpha}_t}\, x_0$: media — versión escalada de los datos originales.
    \item $(1 - \bar{\alpha}_t) I$: varianza — cantidad total de ruido acumulado.
\end{itemize}

\begin{importantbox}{Importancia de la distribución de salto hacia adelante}
La distribución de salto hacia adelante $q(x_t|x_0)$ permite que, durante el entrenamiento, \textbf{directamente} muestremos $x_t$ en cualquier paso temporal a partir de $x_0$, sin necesidad de ejecutar el proceso de difusión paso a paso. Este es el factor clave para la eficiencia del entrenamiento de DDPM. Usando técnicas de reparametrización:
$$x_t = \sqrt{\bar{\alpha}_t}\, x_0 + \sqrt{1 - \bar{\alpha}_t}\, \epsilon, \quad \epsilon \sim \mathcal{N}(0, I)$$
\end{importantbox}

Esta fórmula establece la relación exacta entre $x_0$, $x_t$ y el ruido $\epsilon$, y se volverá a utilizar repetidamente al derivar el objetivo de predicción de ruido.

\subsection{Resumen de la sección}

El proceso de difusión inyecta ruido en los datos gradualmente mediante distribuciones gaussianas de transición predefinidas. Dado que la superposición de distribuciones gaussianas sigue siendo gaussiana, podemos deducir una distribución de salto hacia adelante cerrada $q(x_t|x_0)$, lo que nos permite muestrear $x_t$ directamente desde $x_0$ en un solo paso durante el entrenamiento, sin necesidad de iterar.

%% ========== Sección 3: Descomposición del ELBO y objetivo de entrenamiento ==========


\section{Descomposición del ELBO y objetivo de entrenamiento}

\subsection{De la verosimilitud logarítmica al ELBO}

Similar a los VAE, el objetivo de entrenamiento de los modelos difusos es maximizar la verosimilitud logarítmica de los datos $\log p_\theta(x_0)$. Dado que optimizarlo directamente no es factible, pasamos a optimizar su cota inferior: la cota de evidencia (ELBO):

$$\log p_\theta(x_0) \geq \text{ELBO} = \mathbb{E}_{q(x_{1:T}|x_0)}\left[\log \frac{p_\theta(x_{0:T})}{q(x_{1:T}|x_0)}\right]$$

Equivalentemente, durante el entrenamiento minimizamos el ELBO negativo:

$$L = -\text{ELBO}$$

\begin{knowledgebox}{Relación entre ELBO y VAE}
La derivación del ELBO es idéntica a la de los VAE: se introduce una posterior variacional $q(x_{1:T}|x_0)$ y se utiliza la desigualdad de Jensen para obtener la cota inferior. La diferencia radica en que, para los modelos difusos, $q$ es un proceso hacia adelante predefinido (no aprendido), por lo que el gap del ELBO está determinado por el diseño del proceso hacia adelante.
\end{knowledgebox}

\subsection{Descomposición término a término del ELBO}

Aprovechando la propiedad de Markov y las reglas de Bayes, el ELBO puede descomponerse en la suma de $T+1$ términos. Concretamente, el ELBO negativo se puede expresar como:

$$L = \underbrace{D_{\text{KL}}\big(q(x_T|x_0) \,\|\, p(x_T)\big)}_{L_T\text{(término de ajuste al prior)}} + \sum_{t=2}^{T} \underbrace{D_{\text{KL}}\big(q(x_{t-1}|x_t, x_0) \,\|\, p_\theta(x_{t-1}|x_t)\big)}_{L_{t-1}\text{(término de ajuste al desruido)}} + \underbrace{\big(-\log p_\theta(x_0|x_1)\big)}_{L_0\text{(término de reconstrucción)}}$$

Significado de cada término:
\begin{itemize}
    \item $L_T$ (término de ajuste al prior): divergencia KL entre la distribución terminal del proceso hacia adelante $q(x_T|x_0)$ y el prior $p(x_T) = \mathcal{N}(0,I)$. Dado que $q$ está predefinido y cuando $T$ es suficientemente grande $q(x_T|x_0) \approx \mathcal{N}(0,I)$, este término se aproxima a 0 y \textbf{no requiere optimización}.
    \item $L_{t-1}$ (término de ajuste al desruido): este es el término central; exige que la transferencia inversa aprendida $p_\theta(x_{t-1}|x_t)$ coincida con la transferencia posterior real $q(x_{t-1}|x_t, x_0)$.
    \item $L_0$ (término de reconstrucción): verosimilitud logarítmica negativa del último paso para reconstruir $x_0$ a partir de $x_1$.
\end{itemize}

\begin{importantbox}{El núcleo del entrenamiento: el término de ajuste al desruido $L_{t-1}$}
La meta clave del entrenamiento es minimizar el término de ajuste al desruido $L_{t-1}$, haciendo que la distribución de transferencia inversa $p_\theta(x_{t-1}|x_t)$ se aproxime lo más posible a la transferencia posterior real $q(x_{t-1}|x_t, x_0)$. Esto requiere primero derivar una expresión cerrada para $q(x_{t-1}|x_t, x_0)$.
\end{importantbox}

\subsection{Derivación de la distribución posterior real de transferencia}

Utilizando las reglas de Bayes:

$$q(x_{t-1}|x_t, x_0) = \frac{q(x_t|x_{t-1}, x_0) \cdot q(x_{t-1}|x_0)}{q(x_t|x_0)}$$

Dado que las tres distribuciones $q(x_t|x_{t-1})$, $q(x_{t-1}|x_0)$ y $q(x_t|x_0)$ son distribuciones gaussianas, aprovechando las reglas de multiplicación y división de distribuciones gaussianas, se puede demostrar que $q(x_{t-1}|x_t, x_0)$ también es una distribución gaussiana:

$$q(x_{t-1}|x_t, x_0) = \mathcal{N}\big(x_{t-1};\, \tilde{\mu}_t(x_t, x_0),\; \tilde{\beta}_t I\big)$$

Donde la varianza y la media son respectivamente:

$$\tilde{\beta}_t = \frac{(1 - \bar{\alpha}_{t-1}) \cdot \beta_t}{1 - \bar{\alpha}_t}$$

$$\tilde{\mu}_t(x_t, x_0) = \frac{\sqrt{\bar{\alpha}_{t-1}}\, \beta_t}{1 - \bar{\alpha}_t}\, x_0 + \frac{\sqrt{\alpha_t}\, (1 - \bar{\alpha}_{t-1})}{1 - \bar{\alpha}_t}\, x_t$$

\begin{warningbox}{Dependencia posterior de $x_0$}
Obsérvese que la media $\tilde{\mu}_t$ de $q(x_{t-1}|x_t, x_0)$ depende tanto de $x_t$ como de $x_0$. Sin embargo, en el proceso inverso (proceso generativo), \textbf{no} tenemos información sobre $x_0$; este es precisamente el motivo por el cual necesitamos una red neuronal para predecirlo. La varianza $\tilde{\beta}_t$ solo depende de los parámetros del esquema de ruido, no de $x_0$, por lo que no requiere aprendizaje.
\end{warningbox}

\subsection{Resumen de la sección}

La descomposición del ELBO transforma el objetivo de entrenamiento en minimizar la divergencia KL entre la distribución de transferencia inversa y la distribución posterior real de transferencia en cada paso temporal. La transferencia posterior real $q(x_{t-1}|x_t, x_0)$ es una distribución gaussiana cuya varianza se puede calcular directamente y cuya media es una combinación lineal de $x_t$ y $x_0$.

%% ========== Sección 4: Parametrización del proceso inverso ==========


\section{Parametrización del proceso inverso}

\subsection{Modelado de la distribución de transferencia inversa}

Dado que el posterior real $q(x_{t-1}|x_t, x_0)$ es una distribución gaussiana, naturalmente modelamos también la transferencia inversa $p_\theta(x_{t-1}|x_t)$ como una distribución gaussiana:

$$p_\theta(x_{t-1}|x_t) = \mathcal{N}\big(x_{t-1};\, \mu_\theta(x_t, t),\; \sigma_t^2 I\big)$$

Donde:
\begin{itemize}
    \item $\mu_\theta(x_t, t)$: la media predicha por una red neuronal, tomando como entrada $x_t$ y el paso de tiempo $t$
    \item $\sigma_t^2$: la varianza, fijada en $\tilde{\beta}_t$ (igual que en el posterior real), \textbf{sin necesidad de aprenderla}
\end{itemize}

\begin{knowledgebox}{¿Por qué no es necesario aprender la varianza?}
Dado que queremos minimizar $D_{\text{KL}}(q(x_{t-1}|x_t,x_0) \| p_\theta(x_{t-1}|x_t))$, y la divergencia de KL entre dos distribuciones gaussianas se simplifica a la norma L2 de la diferencia de medias cuando las varianzas son iguales. Por lo tanto, fijar la varianza de $p_\theta$ en $\tilde{\beta}_t$, igual que en $q$, hace que el objetivo de entrenamiento se reduzca únicamente a ajustar la media. El artículo Improved DDPM explora esquemas para aprender la varianza, pero en DDPM básico es suficiente con fijarla.
\end{knowledgebox}

\subsection{Objetivo de entrenamiento simplificado}

Sustituyendo la forma gaussiana en la divergencia de KL, el término de coincidencia de des ruido $L_{t-1}$ se simplifica a:

$$L_{t-1} = \frac{1}{2\sigma_t^2} \left\| \tilde{\mu}_t(x_t, x_0) - \mu_\theta(x_t, t) \right\|^2 + C$$

Donde $C$ es una constante independiente de $\theta$. Por lo tanto, el objetivo de entrenamiento consiste en que la media predicha por la red neuronal $\mu_\theta(x_t, t)$ se acerque lo más posible a la media del posterior real $\tilde{\mu}_t(x_t, x_0)$.

\begin{importantbox}{Intuición del entrenamiento}
El proceso de entrenamiento puede describirse brevemente como: (1) muestrear un dato limpio $x_0$; (2) muestrear el paso de tiempo $t$; (3) muestrear datos con ruido $x_t$ mediante la distribución de salto hacia adelante; (4) predecir la media $\mu_\theta(x_t, t)$ con una red neuronal; (5) calcular la pérdida L2 entre la media predicha y la real, y realizar el retroceso.
\end{importantbox}

En el entrenamiento práctico, el coeficiente $\frac{1}{2\sigma_t^2}$ equivale a un término de peso dependiente del tiempo $t$. Ho et al. (2020) descubrieron que, en la práctica, \textbf{omitir este peso} (es decir, usar pesos uniformes para todos los pasos de tiempo) funciona muy bien e incluso mejor. Este es un hallazgo empírico importante del artículo sobre DDPM.

\subsection{Resumen de la sección}

Al modelar la transferencia inversa como una distribución gaussiana y fijar la varianza, el objetivo de entrenamiento se simplifica a un problema de regresión L2 de medias. El término de peso suele omitirse en la práctica para simplificar el entrenamiento.

%% ========== Sección 5: Tres objetivos de predicción ==========


\section{Tres objetivos de predicción equivalentes}

Dado que la media posterior real $\tilde{\mu}_t(x_t, x_0)$ es una función lineal de $x_t$ y $x_0$, y existe una relación determinista entre $x_0$, $x_t$ y el ruido $\epsilon$, podemos reparametrizar la predicción de la media en tres objetivos de predicción equivalentes.

\subsection{Enfoque 1: Predicción de la media (Mean Prediction)}

La forma más directa consiste en que la red neuronal prediga directamente la media de la transición inversa:

$$\mu_\theta(x_t, t) \approx \tilde{\mu}_t(x_t, x_0)$$

La red acepta $x_t$ y $t$ como entrada y produce un vector de medias con la misma dimensión que $x_t$.

La función de pérdida de entrenamiento es:

$$L_{\text{mean}} = \mathbb{E}_{t, x_0, x_t}\left[\left\| \tilde{\mu}_t(x_t, x_0) - \mu_\theta(x_t, t) \right\|^2\right]$$

\begin{knowledgebox}{Características de la predicción de la media}
La predicción de la media es la forma de parametrización más natural; corresponde directamente al objetivo de optimización derivado del ELBO. Sin embargo, la escala de la media $\tilde{\mu}_t$ varía con el paso temporal $t$ (ya que es una combinación ponderada de $x_t$ y $x_0$, donde los pesos dependen de $t$), lo que puede provocar un entrenamiento poco estable.
\end{knowledgebox}

\subsection{Enfoque 2: Predicción de la muestra limpia (Clean Sample Prediction / $x_0$-Prediction)}

Dado que $\tilde{\mu}_t(x_t, x_0)$ es una función lineal de $x_t$ y $x_0$, y $x_t$ se conoce durante la inferencia, podemos pedirle a la red que prediga $x_0$ en lugar de predecir directamente la media:

$$\hat{x}_0 = f_\theta(x_t, t)$$

A continuación, calculamos la media mediante la fórmula:

$$\mu_\theta(x_t, t) = \frac{\sqrt{\bar{\alpha}_{t-1}}\, \beta_t}{1 - \bar{\alpha}_t}\, f_\theta(x_t, t) + \frac{\sqrt{\alpha_t}\, (1 - \bar{\alpha}_{t-1})}{1 - \bar{\alpha}_t}\, x_t$$

La función de pérdida de entrenamiento equivalente (ignorando los términos de ponderación) es:

$$L_{x_0} = \mathbb{E}_{t, x_0, \epsilon}\left[\left\| x_0 - f_\theta(x_t, t) \right\|^2\right]$$

\begin{importantbox}{Perspectiva única de la predicción $x_0$}
La predicción $x_0$ ofrece una intuición de muestreo atractiva: en cada paso intermedio $t$ del proceso inverso, la red \textbf{predice el valor esperado de la salida final $x_0$}, y luego decide cómo muestrear el siguiente paso $x_{t-1}$ basándose en ello. Esto significa que el proceso de muestreo de los modelos de difusión es un proceso de \textbf{refinamiento iterativo} (iterative refinement): en cada paso se predice el resultado final y, a continuación, la predicción se corrige gradualmente en lugar de simplemente generar una salida directa.
\end{importantbox}

\begin{warningbox}{La predicción $x_0$ no es generación en un solo paso}
Aunque en cada paso se predice $x_0$, no utilizamos directamente esta predicción como salida final. En su lugar, calculamos la media de la transición inversa con $\hat{x}_0$ y luego muestreamos $x_{t-1}$; posteriormente, predecimos nuevamente $\hat{x}_0$ sobre $x_{t-1}$, iterando así hasta llegar a $x_0$. La predicción de $\hat{x}_0$ en cada paso se vuelve gradualmente más precisa.
\end{warningbox}

\subsection{Enfoque 3: Predicción del ruido (Noise Prediction / $\epsilon$-Prediction)}

Esta es la forma de parametrización adoptada en el artículo DDPM y también es la más utilizada en la práctica.

Recordemos la reparametrización de la distribución de salto hacia adelante:

$$x_t = \sqrt{\bar{\alpha}_t}\, x_0 + \sqrt{1 - \bar{\alpha}_t}\, \epsilon, \quad \epsilon \sim \mathcal{N}(0, I)$$

De aquí podemos expresar $x_0$ como una función de $x_t$ y $\epsilon$:

$$x_0 = \frac{x_t - \sqrt{1 - \bar{\alpha}_t}\, \epsilon}{\sqrt{\bar{\alpha}_t}}$$

Sustituyendo esto en la expresión de $\tilde{\mu}_t$, obtenemos una media expresada en términos de $x_t$ y $\epsilon$:

$$\tilde{\mu}_t(x_t, \epsilon) = \frac{1}{\sqrt{\alpha_t}}\left(x_t - \frac{\beta_t}{\sqrt{1 - \bar{\alpha}_t}}\, \epsilon\right)$$

Por lo tanto, podemos pedirle a la red que prediga el ruido $\epsilon$ en lugar de la media o $x_0$:

$$\hat{\epsilon} = \epsilon_\theta(x_t, t)$$

A continuación, calculamos la media mediante la fórmula:

$$\mu_\theta(x_t, t) = \frac{1}{\sqrt{\alpha_t}}\left(x_t - \frac{\beta_t}{\sqrt{1 - \bar{\alpha}_t}}\, \epsilon_\theta(x_t, t)\right)$$

\begin{importantbox}{Objetivo de entrenamiento simplificado de DDPM}
El hallazgo central de Ho et al. (2020): tras eliminar los términos de ponderación, el objetivo de entrenamiento de la predicción del ruido se reduce a una forma extremadamente simple:
$$L_{\text{simple}} = \mathbb{E}_{t, x_0, \epsilon}\left[\left\| \epsilon - \epsilon_\theta(x_t, t) \right\|^2\right]$$
donde $t \sim \text{Uniform}(\{1, \dots, T\})$, $x_0 \sim q(x_0)$, $\epsilon \sim \mathcal{N}(0, I)$ y $x_t = \sqrt{\bar{\alpha}_t}\, x_0 + \sqrt{1 - \bar{\alpha}_t}\, \epsilon$.

Esto es: \textbf{dada una imagen con ruido, predecir el ruido añadido, minimizando el error L2}.
\end{importantbox}

\subsection{Equivalencia y selección entre las tres parametrizaciones}

Las tres parametrizaciones son matemáticamente completamente equivalentes, ya que se pueden convertir unas a otras:

\begin{itemize}
    \item De $\hat{x}_0$ se puede calcular $\hat{\epsilon} = \frac{x_t - \sqrt{\bar{\alpha}_t}\, \hat{x}_0}{\sqrt{1-\bar{\alpha}_t}}$
    \item De $\hat{\epsilon}$ se puede calcular $\hat{x}_0 = \frac{x_t - \sqrt{1-\bar{\alpha}_t}\, \hat{\epsilon}}{\sqrt{\bar{\alpha}_t}}$
    \item De $\hat{x}_0$ o $\hat{\epsilon}$ se puede calcular $\hat{\mu}$
\end{itemize}

\begin{table}[H]
\centering
\caption{Comparación de los tres objetivos de predicción}
\begin{tabular}{llll}
\toprule
\textbf{Forma de parametrización} & \textbf{Objetivo de predicción} & \textbf{Función de pérdida} & \textbf{Observaciones} \\
\midrule
Predicción de la media & $\tilde{\mu}_t$ & $\|\tilde{\mu}_t - \mu_\theta\|^2$ & Más natural pero depende de la escala en $t$ \\
Predicción $x_0$ & $x_0$ & $\|x_0 - f_\theta(x_t,t)\|^2$ & Intuición de refinamiento iterativo \\
Predicción $\epsilon$ & $\epsilon$ & $\|\epsilon - \epsilon_\theta(x_t,t)\|^2$ & Predeterminado en DDPM, más estable \\
\bottomrule
\end{tabular}
\end{table}

\begin{importantbox}{¿Por qué la predicción del ruido es la más común?}
En la práctica, la predicción del ruido ($\epsilon$-prediction) es la forma de parametrización ampliamente utilizada; las razones principales son la \textbf{estabilidad del entrenamiento}:
\begin{enumerate}
    \item \textbf{Normalización del valor objetivo}: el ruido $\epsilon \sim \mathcal{N}(0,I)$ es una muestra de una distribución normal estándar, que posee naturalmente características de escala óptimas con media cero y varianza unitaria. En cambio, la escala de $x_0$ depende del conjunto de datos, y la escala de $\tilde{\mu}_t$ varía con $t$.
    \item \textbf{Consistencia en la escala del gradiente}: dado que el objetivo del ruido tiene las mismas características estadísticas en todos los pasos temporales, la escala de los gradientes es más consistente entre diferentes pasos temporales, lo cual favorece la estabilidad del entrenamiento.
    \item \textbf{Sin normalización adicional}: para los datos $x_0$, generalmente se requiere normalizarlos a un rango específico; mientras que el ruido $\epsilon$ ya está normalizado por sí mismo.
\end{enumerate}
\end{importantbox}

\begin{warningbox}{Diferentes parametrizaciones son adecuadas para diferentes escenarios}
Aunque la predicción del ruido ofrece los mejores resultados en la mayoría de los escenarios, la predicción $x_0$ puede ser más ventajosa en aplicaciones específicas, por ejemplo: edición de imágenes (necesita estimar directamente la imagen limpia), generación de video (necesita mantener la consistencia temporal), entre otras. La parametrización v-prediction posterior combina las ventajas de ambas. Al elegir la forma de parametrización, se deben considerar los requisitos específicos de la tarea.
\end{warningbox}

\subsection{Resumen de la sección}

Los tres objetivos de predicción son matemáticamente equivalentes, pero muestran comportamientos distintos durante el entrenamiento real. La predicción del ruido se ha convertido en la opción por defecto debido a sus buenas características de normalización y estabilidad del entrenamiento. La predicción de muestras limpias ofrece una perspectiva intuitiva de refinamiento iterativo. Las tres formas son intercambiables; en implementaciones prácticas, esta relación de conversión suele aprovecharse durante la inferencia.

%% ========== Sección 6: Algoritmo de entrenamiento ==========
\section{Algoritmo de entrenamiento de DDPM}

\subsection{Flujo completo de entrenamiento}

Combinando las deducciones anteriores, el algoritmo de entrenamiento de DDPM puede escribirse en una forma extremadamente concisa:

\begin{enumerate}
    \item \textbf{Repetir} los siguientes pasos hasta la convergencia:
    \begin{enumerate}
        \item Muestrear un punto de datos limpio del conjunto de entrenamiento: $x_0 \sim q(x_0)$
        \item Muestrear uniformemente y aleatoriamente un paso temporal: $t \sim \text{Uniform}(\{1, 2, \dots, T\})$
        \item Muestrear ruido gaussiano estándar: $\epsilon \sim \mathcal{N}(0, I)$
        \item Calcular los datos con ruido: $x_t = \sqrt{\bar{\alpha}_t}\, x_0 + \sqrt{1 - \bar{\alpha}_t}\, \epsilon$
        \item Calcular la pérdida y actualizar los parámetros: $\nabla_\theta \left\| \epsilon - \epsilon_\theta(x_t, t) \right\|^2$
    \end{enumerate}
\end{enumerate}

\begin{knowledgebox}{Concisión del entrenamiento}
El proceso completo de entrenamiento se puede implementar con solo unas líneas de código en PyTorch. No existe la inestabilidad del entrenamiento adversarial (como en los GAN), ni la complejidad de la inferencia variacional (como el reparametrizado y el equilibrio del término KL en los VAE); es simplemente un problema de regresión L2. Esta concisión es una de las razones principales por las que los modelos de difusión son ampliamente adoptados.
\end{knowledgebox}

\subsection{Implementación pseudocódigo del entrenamiento}

A continuación, se presenta una implementación simplificada con estilo PyTorch:

\begin{lstlisting}[caption={Bucle central del entrenamiento de DDPM}, label=lst:train]
# alpha_bar: producto acumulado precalculado, forma [T]
# model: red neuronal de predicción de ruido epsilon_theta(x_t, t)

for x0 in dataloader:
    t = torch.randint(1, T+1, (batch_size,))
    epsilon = torch.randn_like(x0)

    # Salto hacia adelante: muestrear x_t directamente desde x_0
    sqrt_alpha_bar = alpha_bar[t].sqrt()
    sqrt_one_minus = (1 - alpha_bar[t]).sqrt()
    x_t = sqrt_alpha_bar * x0 + sqrt_one_minus * epsilon

    # Predecir el ruido y calcular la pérdida
    epsilon_pred = model(x_t, t)
    loss = F.mse_loss(epsilon_pred, epsilon)

    optimizer.zero_grad()
    loss.backward()
    optimizer.step()
\end{lstlisting}

\begin{warningbox}{Precálculo de $\bar{\alpha}_t$}
En la implementación, $\bar{\alpha}_t = \prod_{s=1}^{t}(1-\beta_s)$ debe \textbf{precalcularse} y almacenarse como un vector de longitud $T$ antes del inicio del entrenamiento. En cada paso de entrenamiento, se utiliza el valor de $\bar{\alpha}_t$ correspondiente al índice $t$ muestreado aleatoriamente. Tenga en cuenta manejar correctamente la situación donde diferentes muestras en la dimensión batch corresponden a diferentes pasos temporales.
\end{warningbox}

\subsection{Discusión sobre los términos de peso}

Según la deducción estricta del ELBO, el término de coincidencia de desruido $L_{t-1}$ tiene un coeficiente dependiente de $t$, $\frac{1}{2\sigma_t^2}$. Si conservamos este peso, las pérdidas de diferentes pasos temporales se asignarían diferentes importancias.

Ho et al. (2020) descubrieron que en la práctica \textbf{omitir los pesos} (ponderar uniformemente para todos los $t$) produce mejores resultados. Intuitivamente, esto se debe a:

\begin{itemize}
    \item Cuando el entrenamiento es lo suficientemente prolongado y $t$ se muestrea uniformemente, cada paso temporal se entrena adecuadamente
    \item Los pesos uniformes evitan que algunos pasos temporales sean sobreemphasizados mientras otros son ignorados
    \item Experimentalmente, los puntajes FID con pesos uniformes suelen ser mejores
\end{itemize}

\begin{knowledgebox}{Teoría y práctica de los términos de peso}
Aunque teóricamente los términos de peso corresponden a la optimización correcta del ELBO, el objetivo sin pesos puede considerarse una estrategia de \textbf{reponderación}; en realidad, tiende más a mejorar el rendimiento a niveles altos de ruido. Investigaciones posteriores (como P2 weighting, Min-SNR weighting, etc.) exploraron diversas estrategias de peso adaptativas, intentando encontrar un mejor equilibrio entre la corrección teórica y los resultados prácticos.
\end{knowledgebox}

\subsection{Resumen de la sección}

El algoritmo de entrenamiento de DDPM es extremadamente conciso: muestrear datos, muestrear paso temporal, muestrear ruido, calcular datos con ruido, predecir ruido y minimizar la pérdida L2. El proceso completo se puede implementar con menos de 10 líneas de código en PyTorch. Los términos de peso generalmente se omiten en la práctica.

%% ========== Sección siete: Algoritmo de muestreo ==========


````
\section{Algoritmo de muestreo DDPM}

\subsection{Muestreo del proceso inverso}

Una vez completada la fase de entrenamiento, el proceso para generar nuevos muestras (proceso inverso / proceso de muestreo) es el siguiente:

\begin{enumerate}
    \item Muestrear ruido inicial desde una distribución gaussiana estándar: $x_T \sim \mathcal{N}(0, I)$
    \item Para $t = T, T-1, \dots, 1$:
    \begin{enumerate}
        \item Predecir el ruido con la red: $\hat{\epsilon} = \epsilon_\theta(x_t, t)$
        \item Calcular la media: $\mu_\theta(x_t, t) = \frac{1}{\sqrt{\alpha_t}}\left(x_t - \frac{\beta_t}{\sqrt{1-\bar{\alpha}_t}}\, \hat{\epsilon}\right)$
        \item Si $t > 1$, muestrear $z \sim \mathcal{N}(0, I)$; si $t = 1$, establecer $z = 0$
        \item $x_{t-1} = \mu_\theta(x_t, t) + \sigma_t\, z$
    \end{enumerate}
    \item Retornar $x_0$
\end{enumerate}

\begin{importantbox}{Intuición del muestreo: eliminación progresiva de ruido}
El proceso de muestreo puede entenderse como: comenzar con un bloque de puro ruido y, en cada paso, permitir que la red «mire» la imagen actual para predecir cuánto ruido contiene, luego restar el ruido predicho (agregando una pequeña cantidad de nuevo ruido al mismo tiempo para mantener la aleatoriedad), iterando así hasta obtener una imagen limpia. El último paso ($t=1$) no agrega ruido, ya que deseamos una salida determinista.
\end{importantbox}

\subsection{Pseudocódigo del muestreo}

\begin{lstlisting}[caption={Algoritmo de muestreo DDPM}, label=lst:sample]
def ddpm_sample(model, T, alpha, alpha_bar, beta):
    x = torch.randn(shape)  # x_T ~ N(0, I)

    for t in range(T, 0, -1):
        epsilon_pred = model(x, t)

        # Compute mean
        mu = (1 / alpha[t].sqrt()) * (
            x - (beta[t] / (1 - alpha_bar[t]).sqrt()) * epsilon_pred
        )

        if t > 1:
            sigma = beta[t].sqrt()  # or tilde_beta[t].sqrt()
            z = torch.randn_like(x)
            x = mu + sigma * z
        else:
            x = mu  # No noise at final step

    return x
\end{lstlisting}

\begin{warningbox}{La velocidad de muestreo es el principal cuello de botella de DDPM}
El muestreo de DDPM requiere iterar $T$ pasos (usualmente $T = 1000$), y cada paso necesita una propagación hacia adelante completa de la red neuronal. Esto implica que generar una imagen requiere 1000 propagaciones hacia adelante, lo cual es aproximadamente tres órdenes de magnitud más lento que las únicas propagaciones hacia adelante de GAN. Trabajos posteriores como DDIM, DPM-Solver y otros se han centrado en reducir el número de pasos de muestreo, logrando una calidad cercana en tan solo 10--50 pasos.
\end{warningbox}

\subsection{Visualización de $\hat{x}_0$ durante el proceso de muestreo}

Al utilizar una red de predicción de ruido, en cada paso intermedio podemos utilizar la relación para calcular la estimación actual de $x_0$:

$$\hat{x}_0^{(t)} = \frac{x_t - \sqrt{1-\bar{\alpha}_t}\, \epsilon_\theta(x_t, t)}{\sqrt{\bar{\alpha}_t}}$$

En las etapas tempranas del proceso de muestreo ($t$ cercano a $T$), $\hat{x}_0^{(t)}$ es muy difuso; a medida que $t$ disminuye, la predicción se vuelve gradualmente más clara. Esta visualización ilustra intuitivamente la característica de «refinamiento iterativo» de los modelos de difusión.

\begin{knowledgebox}{Refinamiento iterativo y Progressive Refinement}
La predicción de $\hat{x}_0$ en cada paso puede entenderse como: la mejor conjetura sobre el resultado final bajo el nivel actual de información. A medida que avanza la eliminación de ruido, aumenta la cantidad de información disponible (y disminuye el ruido), por lo que la predicción se vuelve naturalmente más precisa. Esta característica de *progressive refinement* es muy útil en aplicaciones prácticas —por ejemplo, para guiar la dirección de generación en pasos intermedios (*classifier guidance* / *classifier-free guidance*)—.
\end{knowledgebox}

\subsection{Resumen de la sección}

El proceso de muestreo de DDPM es la «operación inversa» del proceso de entrenamiento: partiendo de puro ruido, se elimina el ruido gradualmente hasta obtener una imagen limpia. En cada paso se predice el ruido, se calcula la media y se agrega ruido aleatorio controlado. La lentitud del muestreo es la principal limitación, pero existen numerosos métodos de aceleración posteriores.

%% ========== Sección 8: Capacidad expresiva de los modelos de difusión ==========
\section{¿Por qué los modelos de difusión pueden modelar distribuciones complejas}

\subsection{Gaussiano de un paso vs. Gaussiano de múltiples pasos}

El conferenciante planteó en clase una pregunta profunda: dado que las transiciones del proceso hacia adelante son gaussianas y la transición del proceso hacia atrás también se modela como gaussiana, ¿la distribución de verosimilitud completa $p_\theta(x_0|x_T)$ también es una distribución gaussiana? Si lo fuera, los modelos de difusión no estarían tan limitados como los VAE.

La respuesta es \textbf{no}.

\begin{importantbox}{La no linealidad introduce complejidad}
Aunque cada transición hacia atrás individual $p_\theta(x_{t-1}|x_t)$ es una distribución gaussiana, la distribución de verosimilitud completa $p_\theta(x_0|x_T) = \int p_\theta(x_0|x_1) p_\theta(x_1|x_2) \cdots p_\theta(x_{T-1}|x_T) \, dx_{1:T-1}$ \textbf{no} es una distribución gaussiana. Lo esencial está en que la \textbf{media} de cada paso gaussiano se predice mediante una red neuronal (una función no lineal).

A diferencia del proceso hacia adelante —donde la media es una función lineal de $x_{t-1}$ ($\sqrt{1-\beta_t}\, x_{t-1}$), por lo que la combinación de múltiples pasos en el proceso hacia adelante sigue siendo gaussiana—, en el proceso hacia atrás la media es una función no lineal de $x_t$ (a través de la red neuronal). Esto permite que la combinación de múltiples pasos exprese cualquier distribución compleja.
\end{importantbox}

\begin{warningbox}{Proceso hacia adelante lineal vs. proceso hacia atrás no lineal}
Esto es fundamental para entender la capacidad de expresión de los modelos de difusión:
\begin{itemize}
    \item \textbf{Proceso hacia adelante}: La media es una función \textbf{lineal} de la entrada $\rightarrow$ la combinación de múltiples pasos sigue siendo gaussiana $\rightarrow$ se puede derivar una distribución de salto en forma cerrada.
    \item \textbf{Proceso hacia atrás}: La media es una función \textbf{no lineal} de la entrada (red neuronal) $\rightarrow$ la combinación de múltiples pasos \textbf{no} es gaussiana $\rightarrow$ la distribución de verosimilitud puede ser arbitrariamente compleja.
\end{itemize}
Los modelos de difusión encuentran un equilibrio inteligente entre ``paso simple en cada etapa de entrenamiento'' y ``fuerte capacidad de expresión global''.
\end{warningbox}

Esto explica también por qué los modelos de difusión requieren múltiples pasos: cuanto más es el número de pasos, más profundas son las transformaciones no lineales y mayor es la capacidad de expresión. Esto es análogo a la razón por la que las redes neuronales profundas tienen mayor capacidad de expresión que las redes superficiales.

\subsection{Resumen de la sección}

Aunque cada transición del modelo de difusión es gaussiana, debido a que el proceso hacia atrás utiliza una red neuronal no lineal para predecir la media, la distribución de verosimilitud global resultante de la combinación de múltiples pasos puede modelar distribuciones arbitrariamente complejas. Esta es la razón fundamental por la cual los modelos de difusión son superiores a los VAE simples.

%% ========== Novena sección: Conexión entre predicción de ruido y emparejamiento de puntuaciones ==========


\section{Relación entre la predicción de ruido y el \textit{score matching}}

\subsection{De la predicción de ruido a la función \textit{score}}

En las conclusiones del curso, el instructor anunció una conexión importante: existe una correspondencia profunda entre las redes de predicción de ruido $\epsilon_\theta(x_t, t)$ y la función \textit{score} en los modelos basados en \textit{score}.

La función \textit{score} se define como el gradiente de la densidad de probabilidad logarítmica respecto a los datos:

$$s_\theta(x_t, t) = \nabla_{x_t} \log p_t(x_t)$$

Se puede demostrar que existe la siguiente relación entre la red de predicción de ruido y la función \textit{score}:

$$\epsilon_\theta(x_t, t) \approx -\sqrt{1 - \bar{\alpha}_t}\, \nabla_{x_t} \log q_t(x_t)$$

\begin{knowledgebox}{Unificación de DDPM y los modelos basados en Score}
Esta conexión revela que dos marcos de modelos generativos que a primera vista parecen distintos —los modelos probabilísticos de difusión por eliminación de ruido (DDPM) y los modelos generativos basados en score (SGM)— son, en esencia, equivalentes:
\begin{itemize}
    \item DDPM aprende a predecir el ruido $\epsilon$
    \item SGM aprende a predecir la función \textit{score} $\nabla_x \log p(x)$
    \item La única diferencia entre ambos es un factor de escala relacionado con el paso temporal
\end{itemize}
Esta equivalencia fue descrita en una descripción continua unificada dentro del marco Score SDE por Song et al. (2021).
\end{knowledgebox}

\subsection{Comprensión intuitiva}

La predicción de ruido puede entenderse como responder a la siguiente pregunta: ``¿cuánto ruido hay en $x_t$?'' Por otro lado, la función \textit{score} responde: ``¿en qué dirección debería moverse $x_t$ para aumentar su densidad de probabilidad?'' En el contexto del proceso de difusión, estas dos preguntas son esencialmente equivalentes: la dirección en la que se elimina el ruido es exactamente la dirección en la que aumenta la densidad de probabilidad.

\subsection{Resumen de la sección}

Existe una correspondencia matemática precisa entre la predicción de ruido y la función \textit{score}, lo cual establece un puente entre los modelos DDPM y los basados en score, sentando las bases para el marco unificado posterior de Score SDE.

%% ========== Síntesis y ampliación ==========


\section{Síntesis y ampliación}

\subsection{Repaso de los puntos clave}

En esta clase se derivó detalladamente el algoritmo de entrenamiento y muestreo de DDPM; su contenido central puede resumirse en lo siguiente:

\begin{enumerate}
    \item El \textbf{proceso hacia adelante} es un proceso predefinido de inyección de ruido gaussiano, con una distribución de salto de forma cerrada $q(x_t|x_0)$, que permite el muestreo en un solo paso.
    \item La \textbf{descomposición del ELBO} transforma el objetivo de entrenamiento para que coincida con la distribución posterior real de transición en cada paso temporal.
    \item La \textbf{posterior real} $q(x_{t-1}|x_t, x_0)$ es una distribución gaussiana cuya varianza puede calcularse directamente y cuyo valor medio depende de $x_0$.
    \item Tres \textbf{parametrizaciones}: predicción del valor medio, predicción de $x_0$ y predicción de ruido son matemáticamente equivalentes; en la práctica, la \textbf{predicción de ruido} es la más estable.
    \item \textbf{Objetivo simplificado de entrenamiento}: $L_{\text{simple}} = \mathbb{E}\|\epsilon - \epsilon_\theta(x_t, t)\|^2$, elegante y eficiente.
    \item La \textbf{capacidad de expresión} proviene de la no linealidad introducida por la red neuronal en el proceso inverso: gaussiano paso a paso pero globalmente no gaussiano.
    \item Existe una conexión profunda entre la \textbf{predicción de ruido} y los modelos basados en puntuación (score-based), sentando las bases para un marco teórico unificado.
\end{enumerate}

\subsection{Anticipación de las próximas clases}

El ponente mencionará que se discutirán los siguientes temas:
\begin{itemize}
    \item La conexión formal entre predicción de ruido y score matching.
    \item El marco Score SDE: extensión del tiempo discreto de difusión al continuo.
    \item Diferentes definiciones del proceso hacia adelante.
    \item Métodos de muestreo más rápidos (como DDIM).
\end{itemize}

\subsection{Lecturas adicionales}

\begin{itemize}
    \item \textbf{Artículo original de DDPM}: Ho, Jain, Abbeel. ``Denoising Diffusion Probabilistic Models.'' NeurIPS 2020.
    \item \textbf{DDPM mejorado}: Nichol, Dhariwal. ``Improved Denoising Diffusion Probabilistic Models.'' ICML 2021. Explora estrategias para aprender la varianza.
    \item \textbf{Score SDE}: Song et al. ``Score-Based Generative Modeling through Stochastic Differential Equations.'' ICLR 2021. Unifica los modelos DDPM y score-based.
    \item \textbf{DDIM}: Song, Meng, Ermon. ``Denoising Diffusion Implicit Models.'' ICLR 2021. Propone un método de muestreo determinista que reduce drásticamente el número de pasos de muestreo.
    \item \textbf{P2 Weighting}: Choi et al. ``Perception Prioritized Training of Diffusion Models.'' CVPR 2022. Explora estrategias adaptativas de ponderación del paso temporal.
    \item \textbf{v-prediction}: Salimans, Ho. ``Progressive Distillation for Fast Sampling of Diffusion Models.'' ICLR 2022. Propone la parametrización v-prediction, que combina las ventajas de la predicción de $x_0$ y la predicción de $\epsilon$.
\end{itemize}

\end{document}
